\documentclass[12pt,a4paper]{report}

% Packages
\usepackage[margin=1in]{geometry}
\usepackage{graphicx}
\usepackage{times}
\usepackage{setspace}
\usepackage{titlesec}
\usepackage{tocloft}
\usepackage{hyperref}
\usepackage{booktabs}
\usepackage{array}
\usepackage{enumitem}
\usepackage{fancyhdr}
\usepackage{caption}
\usepackage{xcolor}
\usepackage{listings}
\usepackage{amsmath}
\usepackage{float}
\usepackage{parskip}

% Hyperref Setup
\hypersetup{
    colorlinks=true,
    linkcolor=black,
    urlcolor=blue,
    citecolor=black
}

% Section Title Formatting
\titleformat{\chapter}[display]
  {\normalfont\Large\bfseries\centering}
  {Chapter \thechapter}{10pt}{\Large}
\titlespacing*{\chapter}{0pt}{20pt}{20pt}

\titleformat{\section}{\normalfont\large\bfseries}{\thesection}{1em}{}
\titleformat{\subsection}{\normalfont\normalsize\bfseries}{\thesubsection}{1em}{}

% Header / Footer
\pagestyle{fancy}
\fancyhf{}
\rhead{\small Resume Gap Analysis using RAG}
\lhead{\small Dr. D. Y. Patil Institute of Technology}
\cfoot{\thepage}
\renewcommand{\headrulewidth}{0.4pt}

% Code Listing Style
\lstset{
  basicstyle=\ttfamily\small,
  breaklines=true,
  frame=single,
  backgroundcolor=\color{gray!10},
  keywordstyle=\color{blue},
  commentstyle=\color{green!60!black},
  stringstyle=\color{red!70!black}
}

% TOC Formatting
\renewcommand{\cftchapfont}{\bfseries}
\renewcommand{\cftsecfont}{\normalfont}
\setlength{\cftbeforesecskip}{2pt}

\begin{document}

% TITLE PAGE
\begin{titlepage}
\begin{center}

% Institution logo - replace with actual logo file if available
\includegraphics[width=0.30\textwidth]{dypdpu_logo.png}\\[0.3cm]
{\normalsize \textcolor{red!70!black}{Dr. D. Y. Patil Institute of Technology}}\\[0.5cm]
% \rule{\linewidth}{0.5pt}\\[0.3cm]
{\normalsize A}\\[0.1cm]
{\normalsize Data Science and Big Data Analytics}\\
{\normalsize Mini Project Report submitted to Savitribai Phule Pune University, Pune}\\[0.5cm]
\includegraphics[width=0.22\textwidth]{Reportsppu.png}\\[0.3cm]

% \rule{\linewidth}{0.5pt}\\[1cm]

{\Huge \textbf{RESUME GAP ANALYSIS}}\\[0.4cm]
{\Huge \textbf{USING RAG AND FLASK}}\\[1cm]

% University seal placeholder
% \includegraphics[width=0.18\textwidth]{sppu_seal.png}\\[0.8cm]

{\normalsize In partial fulfilment for the awards of Degree of Engineering in}\\
{\normalsize Computer Engineering}\\[1cm]

{\normalsize \textbf{Submitted by}}\\[0.4cm]
\begin{tabular}{ll}
\textbf{Mr. Parth Pakhare} & \textbf{Exam Seat No. } \\
\textbf{Miss. Mansi Gangurde} & \textbf{Exam Seat No. } \\
\textbf{Mr. Aditya Kokate} & \textbf{Exam Seat No. } \\
\end{tabular}\\[1cm]

{\normalsize \textbf{Under the Guidance of}}\\[0.3cm]
{\Large \textbf{Prof. Rubi Mandal}}\\
{\Large \textbf{Assistant Professor}}\\[1cm]

{\normalsize \textbf{Department of Computer Engineering}}\\[0.5cm]
{\normalsize \textbf{2025--26}}

\end{center}
\end{titlepage}

% CERTIFICATE PAGE
\thispagestyle{empty}
\begin{center}
{\Large \textbf{Certificate}}
\end{center}
\vspace{0.5cm}

\noindent This is to certify that,

\vspace{0.5cm}
\begin{tabular}{ll}
Mr. Parth Pakhare & Exam Seat No.  \\
Miss. Mansi Gangurde & Exam Seat No.  \\
Mr. Aditya Kokate & Exam Seat No.  \\
\end{tabular}

\vspace{0.5cm}
\noindent have successfully completed the Mini project entitled
\textbf{``Resume Gap Analysis using RAG''} under my guidance in partial
fulfilment of the requirements for the Third Year of Engineering in Computer
Engineering under the Savitribai Phule Pune University during the academic
year 2025--2026.

\vspace{2cm}
\noindent Date: \ldots\ldots\ldots\ldots\ldots \\[0.3cm]
\noindent Place: \ldots\ldots\ldots\ldots\ldots

\vspace{2cm}
\begin{tabular}{p{5cm} p{5cm}}
\textbf{Prof. Rubi Mandal} & \textbf{Dr. Chaya Ravi Jadhav} \\
Project Guide & HOD \\[2cm]
\multicolumn{2}{l}{\textbf{Dr. Nitin Sherje}} \\
\multicolumn{2}{l}{Principal} \\
\end{tabular}

\newpage

% TABLE OF CONTENTS
\tableofcontents
\newpage

% CHAPTER 1: INTRODUCTION
\chapter{Introduction}

\section{Introduction}

In today's highly competitive job market, the alignment between a candidate's
existing skills and the requirements of a target job role has become more critical
than ever. Recruiters and candidates alike face the challenge of quickly and
accurately identifying the gaps in a resume with respect to a specific career
path. Traditional manual processes are time-consuming, inconsistent, and do not
scale well for modern placement and recruitment needs.

The emergence of \textbf{Retrieval-Augmented Generation (RAG)} -- a technique
that combines the retrieval of relevant information from a knowledge base with
the generative power of large language models (LLMs) -- provides an opportunity
to build intelligent, context-aware applications. When applied to career
development, RAG can be used to semantically compare a candidate's resume
against role-specific skill roadmaps, identify skill gaps, and generate a
personalised learning roadmap.

This project implements a \textbf{Flask-based web application} that performs
intelligent Resume Gap Analysis. The system accepts a candidate's resume either
as pasted plain text or as a text-based PDF upload. It also allows the user to
choose from predefined roles such as Frontend Developer, Backend Developer,
DevOps Engineer, Data Scientist, Full Stack Developer, and Android Developer,
or enter a custom target role. The backend extracts resume text, retrieves
relevant roadmap skill chunks from a local vector store, and uses Gemini through
LangChain to generate a detailed gap analysis report and actionable roadmap.

Key components of the proposed system include:

\begin{enumerate}
  \item \textbf{Resume Ingestion and Parsing:} Extraction of text from pasted
        resume content or uploaded PDF resumes using Python and \texttt{pypdf}.
  \item \textbf{Vector Store and Embeddings:} Converting roadmap skill chunks
        into Gemini embedding vectors and storing them in a local SQLite-based
        vector store for semantic retrieval.
  \item \textbf{Retrieval-Augmented Generation:} Retrieving the most relevant
        roadmap chunks using cosine similarity and passing them to Gemini as
        focused context for analysis.
  \item \textbf{Roadmap Generation:} Producing a structured, role-specific
        learning roadmap that guides the candidate through high-priority,
        medium-priority, and low-priority skill gaps.
  \item \textbf{Flask Web Interface:} A lightweight, responsive HTML and CSS
        interface rendered using Flask and Jinja templates.
\end{enumerate}

\section{Problem Addressed}

The primary problem addressed by this project is the \textbf{inefficient and
subjective identification of skill gaps} between a candidate's resume and a
desired job role. Specific challenges include:

\begin{enumerate}
  \item \textbf{Manual Gap Analysis:} Identifying missing skills requires
        significant human effort and domain expertise. This process is slow and
        prone to inconsistency across different evaluators.
  \item \textbf{Lack of Personalisation:} Generic skill-gap tools provide
        one-size-fits-all feedback without considering the candidate's existing
        background or target role.
  \item \textbf{Absence of Actionable Roadmaps:} Most resume screening tools
        highlight gaps but do not provide a structured, prioritised learning path.
  \item \textbf{Scalability in Recruitment:} Organisations receiving large
        volumes of applications cannot feasibly perform manual gap analysis at
        scale.
  \item \textbf{Dynamic Job Market:} Job requirements evolve rapidly; static
        keyword matching approaches fail to capture semantic equivalence between
        skills.
\end{enumerate}

\section{Methods Used}

Several techniques and frameworks are employed in the project to address these
challenges:

\begin{enumerate}
  \item \textbf{Retrieval-Augmented Generation (RAG):}
  \begin{itemize}
    \item RAG combines a retrieval module with an LLM-based generation module.
    \item In this project, role roadmap chunks are retrieved from a local vector
          store and supplied to Gemini as grounded context.
    \item \textit{Strengths:} Context-aware, improves relevance, and reduces
          purely unguided generation.
    \item \textit{Weaknesses:} Depends on the quality of retrieved roadmap chunks
          and embedding similarity.
  \end{itemize}

  \item \textbf{Gemini Embeddings and Vector Similarity:}
  \begin{itemize}
    \item Gemini embedding models convert roadmap text and resume-role queries
          into high-dimensional vectors.
    \item Cosine similarity is used to retrieve roadmap chunks most relevant to
          the candidate's resume and selected role.
    \item \textit{Strengths:} Captures semantic similarity beyond keyword
          matching.
    \item \textit{Weaknesses:} Requires API availability and correct model
          configuration.
  \end{itemize}

  \item \textbf{SQLite-Based Local Vector Store:}
  \begin{itemize}
    \item A local SQLite database stores roadmap chunks and their embedding
          vectors.
    \item \textit{Strengths:} Simple, portable, serverless, and suitable for a
          mini-project.
    \item \textit{Weaknesses:} Less scalable than specialised vector databases
          for very large knowledge bases.
  \end{itemize}

  \item \textbf{Large Language Model using Gemini:}
  \begin{itemize}
    \item Gemini is used to generate structured JSON output containing match
          score, summary, strengths, matched skills, missing skills, and roadmap.
    \item \textit{Strengths:} Produces readable, detailed, and structured output.
    \item \textit{Weaknesses:} Depends on prompt quality and external API
          availability.
  \end{itemize}

  \item \textbf{Flask Web Framework:}
  \begin{itemize}
    \item Flask exposes the RAG pipeline through server-rendered web pages.
    \item \textit{Strengths:} Simple, flexible, and easy to integrate with Python
          AI libraries.
    \item \textit{Weaknesses:} Requires additional production tooling for
          high-concurrency deployment.
  \end{itemize}
\end{enumerate}

\section{Gaps and Limitations}

Despite recent advances in NLP and RAG-based systems, several gaps remain:

\begin{enumerate}
  \item \textbf{Domain-Specific Knowledge:} Generic LLMs may not always reflect
        niche or emerging technologies, requiring richer and regularly updated
        knowledge bases.
  \item \textbf{Resume Format Variability:} Resumes come in highly varied
        formats, and PDF parsing may lose structural context such as tables,
        columns, or icons.
  \item \textbf{Scanned PDF Limitation:} The current implementation supports
        text-based PDFs. Image-only scanned resumes require OCR, which is not
        included in the present version.
  \item \textbf{Lack of Ground Truth:} There is no universally agreed benchmark
        dataset for resume-to-role gap analysis, making objective evaluation
        difficult.
  \item \textbf{Real-Time Roadmap Updates:} The generated roadmap does not
        dynamically update as the candidate completes learning milestones.
\end{enumerate}

\section{Mini Project Idea}

\textbf{Project Title:} ``Intelligent Resume Gap Analysis and Career Roadmap
Generation using RAG and Flask''

\textbf{Objective:} To develop an AI-powered web application that (1)
semantically analyses a candidate's resume against a specified job role, (2)
identifies specific skill and experience gaps, and (3) generates a structured,
prioritised learning roadmap to help the candidate become job-ready.

\textbf{Proposed Methodology:}

\begin{enumerate}
  \item \textbf{Input Collection:}
        Paste resume text or upload a text-based resume PDF and select a
        predefined or custom target role through the Flask web interface.
  \item \textbf{Resume Processing:}
        Extract text from the resume input and validate minimum content length.
  \item \textbf{Roadmap Vector Retrieval:}
        Embed predefined roadmap skill chunks using Gemini embeddings and store
        them in a local SQLite vector store.
  \item \textbf{RAG Context Selection:}
        Embed the resume-role query and retrieve the most relevant roadmap chunks
        using cosine similarity.
  \item \textbf{Gap Analysis Generation:}
        Pass resume text, target role, and retrieved roadmap context to Gemini
        through a LangChain prompt-model-parser chain.
  \item \textbf{Web Interface:}
        Display the match score, strengths, matched skills, missing skills, and
        learning roadmap through a clean Flask-rendered HTML dashboard.
\end{enumerate}

\section{Expected Outcomes}

\begin{itemize}
  \item Automated identification of skill, tool, and experience gaps between a
        candidate's resume and a target job role.
  \item A structured, prioritised, and time-bound career development roadmap.
  \item Reduced time and effort for candidates and career guidance teams.
  \item A reusable, extensible Flask-based web application deployable for
        institutional career guidance centres.
\end{itemize}

% CHAPTER 2: SYSTEM ARCHITECTURE & METHODOLOGY
\chapter{System Architecture \& Methodology}

\section{System Overview}

The proposed system follows a modular pipeline architecture consisting of five
primary stages: \textbf{Input Processing}, \textbf{Embedding \& Vector Storage},
\textbf{Retrieval}, \textbf{Generation}, and \textbf{Presentation}. Flask serves
as the application layer connecting all modules.

\begin{figure}[H]
  \centering
  % Replace with actual architecture diagram image file
  \includegraphics[width=0.85\textwidth]{system_architecture.png}
  % \fbox{\parbox{0.8\textwidth}{\centering \vspace{2cm}
  % [System Architecture Diagram: Input $\rightarrow$ Parser $\rightarrow$
  % Embeddings $\rightarrow$ SQLite Vector Store $\rightarrow$ Gemini $\rightarrow$ Output]
  % \vspace{2cm}}}
  \caption{High-level system architecture of the Resume Gap Analysis application.}
\end{figure}

\section{Technology Stack}

\begin{table}[H]
\centering
\caption{Technology stack used in the project.}
\begin{tabular}{|p{3.5cm}|p{4cm}|p{5.5cm}|}
\hline
\textbf{Component} & \textbf{Technology} & \textbf{Purpose} \\
\hline
Backend Language & Python & Main programming language for web backend, AI orchestration, PDF parsing, and data handling \\
\hline
Web Framework & Flask & Server-rendered web application, routing, validation, and form handling \\
\hline
Template Engine & Jinja2 & Dynamic rendering of forms, errors, selected roles, and generated results \\
\hline
Frontend & HTML5 and CSS3 & Lightweight user interface without React or frontend JavaScript \\
\hline
PDF Parsing & pypdf & Extract text from uploaded text-based resume PDFs \\
\hline
LLM Orchestration & LangChain & Prompt-template, model, and output-parser chain management \\
\hline
LLM & Gemini via langchain-google-genai & Gap analysis, match scoring, strengths, and roadmap generation \\
\hline
Embeddings & Gemini Embeddings & Generate semantic vectors for roadmap chunks and resume-role queries \\
\hline
Vector Store & SQLite with cosine similarity & Local storage and retrieval of embedded roadmap chunks \\
\hline
Configuration & python-dotenv & Load API keys and model settings from environment variables \\
\hline
\end{tabular}
\end{table}

\section{Dataset and Input Sources}

The system does not rely on a fixed offline resume dataset; instead, it processes
user-supplied input dynamically at runtime:

\begin{itemize}
  \item \textbf{Resume Text:} Entered directly by the user in a textarea. This
        option is useful when the resume is already available as plain text.
  \item \textbf{Resume PDF:} Uploaded by the candidate via the web interface.
        The backend extracts selectable text using \texttt{pypdf}. If both text
        and PDF are provided, the uploaded PDF is used.
  \item \textbf{Target Role:} Selected from predefined roadmap roles or entered
        as a custom target role.
  \item \textbf{Roadmap Knowledge Base:} A built-in collection of role-specific
        roadmap skill categories inspired by roadmap-style learning paths.
\end{itemize}

\section{Data Preprocessing Techniques}

\begin{itemize}
  \item \textbf{PDF Text Extraction:} Raw text is extracted from uploaded PDF
        resumes using \texttt{pypdf}. Image-only scanned PDFs are rejected with
        a validation message.
  \item \textbf{Input Validation:} The backend checks resume length, custom role
        text, PDF file type, PDF size, and missing API key configuration.
  \item \textbf{Roadmap Chunking:} Each predefined roadmap is split line by line
        into skill-category chunks such as databases, APIs, testing, deployment,
        security, and cloud.
  \item \textbf{Embedding Generation:} Roadmap chunks and resume-role queries are
        encoded using Gemini embeddings.
  \item \textbf{SQLite Vector Storage:} Embedding vectors are stored as JSON in a
        local SQLite database and retrieved using cosine similarity.
  \item \textbf{Prompt Construction:} Retrieved roadmap chunks are inserted into
        a structured LangChain prompt along with the resume text and target role.
\end{itemize}

\section{Challenges Faced During Implementation}

\begin{itemize}
  \item \textbf{PDF Parsing Inconsistency:} Highly designed resumes may contain
        tables, columns, or icons that do not extract cleanly. The current system
        handles text-based PDFs and asks users to paste text if extraction fails.
  \item \textbf{Embedding Model Compatibility:} Older embedding model names may
        not be supported by the current Gemini API version. This was resolved by
        configuring the stable Gemini embedding model and rebuilding the vector
        store when the embedding model changes.
  \item \textbf{Vector Database Portability:} Dedicated vector databases can
        require compiled dependencies. A SQLite-based vector store was used to
        keep the mini-project portable and easy to run on Windows.
  \item \textbf{LLM Output Consistency:} The LLM may return extra text unless
        prompted carefully. This was addressed by requesting strict JSON output
        and normalising the parsed result on the backend.
  \item \textbf{API Latency:} Gemini calls and embedding generation may introduce
        response delay, especially during the first run when roadmap vectors are
        created.
\end{itemize}

% CHAPTER 3: EXPLORATORY DATA ANALYSIS & SYSTEM DESIGN
\chapter{Exploratory Analysis \& System Design}

\section{Resume Skill Distribution Analysis}

To understand typical resume profiles, a preliminary conceptual analysis was
considered across common domains such as Software Development, Data Science, and
Cloud Engineering. Skills are commonly grouped into categories such as
\textit{Programming Languages, Frameworks, Databases, Cloud Platforms, Testing,}
and \textit{Soft Skills}. The implemented system follows a similar idea by
organising roadmap knowledge into role-specific skill categories.

\begin{figure}[H]
  \centering
  \includegraphics[width=0.75\textwidth]{skill_distribution.png}
  % \fbox{\parbox{0.75\textwidth}{\centering \vspace{2cm}
  % [Bar Chart: Frequency of skill categories across sample resumes]
  % \vspace{2cm}}}
  \caption{Distribution of skill categories across sample resumes.}
\end{figure}

The analysis highlights that technical resumes often show strong alignment in
programming languages and frameworks, while gaps frequently appear in cloud,
deployment, testing, security, and system design practices.

\section{Semantic Similarity Heatmap}

The system uses cosine similarity between embedding vectors to estimate how
closely a resume-role query matches stored roadmap skill chunks. A heatmap can
be used to visualise such alignment scores, where darker cells indicate higher
semantic similarity.

\begin{figure}[H]
  \centering
  \includegraphics[width=0.75\textwidth]{similarity_heatmap.png}
  % \fbox{\parbox{0.75\textwidth}{\centering \vspace{2cm}
  % [Heatmap: Cosine similarity between resume sections and roadmap requirements]
  % \vspace{2cm}}}
  \caption{Semantic similarity heatmap between resume chunks and role requirements.}
\end{figure}

This similarity-based approach helps identify which roadmap areas are most
relevant to the candidate's resume and target role, allowing Gemini to generate
a focused gap analysis.

\section{RAG Pipeline Flow Diagram}

Figure~3.3 illustrates the end-to-end RAG pipeline implemented in the
application.

\begin{figure}[H]
  \centering
  \includegraphics[width=0.85\textwidth]{rag_pipeline.png}
  % \fbox{\parbox{0.85\textwidth}{\centering \vspace{2.5cm}
  % [Flow Diagram: Resume Text/PDF $\rightarrow$ Parser $\rightarrow$
  % Roadmap Chunks $\rightarrow$ Gemini Embeddings $\rightarrow$
  % SQLite Vector Store $\rightarrow$ Retriever $\rightarrow$
  % Gemini Prompt $\rightarrow$ Gap Report + Roadmap $\rightarrow$ Flask UI]
  % \vspace{2.5cm}}}
  \caption{End-to-end RAG pipeline for resume gap analysis.}
\end{figure}

% CHAPTER 4: IMPLEMENTATION
\chapter{Implementation}

\section{Flask Application Structure}

The application follows a compact Flask project layout:

\begin{lstlisting}[language=bash, caption={Project directory structure.}]
New project/
|-- app.py                         # Flask routes, RAG logic, PDF parsing, Gemini calls
|-- requirements.txt               # Python dependencies
|-- resume_vectors.sqlite3          # Generated local vector store after first RAG run
|-- templates/
|   |-- index.html                 # Form and result page rendered with Jinja
|-- static/
|   |-- styles.css                 # Plain CSS styling
|-- .env                           # Local API key and model settings
|-- .env.example                   # Example environment variables
|-- README.md
\end{lstlisting}

\section{Core RAG Implementation}

The retrieval and generation logic is implemented using LangChain, Gemini
embeddings, a local SQLite vector store, and Gemini chat generation. The key
implementation pattern is shown below:

\begin{lstlisting}[language=Python, caption={RAG pipeline core logic in app.py.}]
def retrieve_roadmap_context(resume, role_context, api_key):
    ensure_roadmap_vector_store(api_key)
    embedder = get_embedding_model(api_key)

    query = (
        f"Target role: {role_context['label']}\n"
        f"Resume excerpt: {resume[:1500]}"
    )
    query_embedding = embedder.embed_query(
        query,
        task_type="RETRIEVAL_QUERY",
    )

    retrieved = query_roadmap_vectors(
        query_embedding,
        role_context.get("role_key", ""),
    )

    chunks = []
    for item in retrieved:
        chunks.append(
            f"[Retrieved: {item['role_label']} / {item['category']}]\n"
            f"{item['text']}"
        )
    return "\n\n".join(chunks)

def analyze_resume(resume, role_context):
    api_key = os.getenv("GEMINI_API_KEY")
    chain = get_analysis_chain(api_key)
    skills_context = retrieve_roadmap_context(resume, role_context, api_key)

    response_text = chain.invoke({
        "role_label": role_context["label"],
        "resume": resume[:6000],
        "skills_context": skills_context,
    })
    return normalize_result(parse_json_response(response_text))
\end{lstlisting}

\section{Flask Route Implementation}

\begin{lstlisting}[language=Python, caption={Main Flask routes in app.py.}]
@app.get("/")
def index():
    return render_template(
        "index.html",
        roles=ROADMAPS,
        selected_role="frontend",
        resume="",
        result=None,
        error="",
    )

@app.get("/analyze")
def analyze_get():
    return redirect(url_for("index"))

@app.post("/analyze")
def analyze():
    pasted_resume = request.form.get("resume", "").strip()
    selected_role = request.form.get("role", "frontend")
    custom_role = request.form.get("custom_role", "").strip()
    resume_pdf = request.files.get("resume_pdf")

    resume = pasted_resume
    if resume_pdf and resume_pdf.filename:
        resume = extract_pdf_text(resume_pdf)

    role_context = get_role_context(selected_role, custom_role)
    result = analyze_resume(resume, role_context)

    return render_template(
        "index.html",
        roles=ROADMAPS,
        selected_role=selected_role,
        custom_role=custom_role,
        role_display=role_context,
        resume=pasted_resume,
        result=result,
        gap_groups=split_gaps(result.get("missingSkills")),
    )
\end{lstlisting}

\section{Prompt Engineering for Gap Analysis}

A structured LangChain prompt template was designed to guide Gemini in producing
consistent, well-formatted JSON output. The prompt includes the candidate's
resume, target role, and retrieved roadmap context. It explicitly requests the
following fields: \textit{matchScore}, \textit{summary}, \textit{strengths},
\textit{matchedSkills}, \textit{missingSkills}, and \textit{roadmap}. Temperature
is set to 0.3 to balance creativity with factual consistency.

% CHAPTER 5: EXPERIMENTAL RESULTS AND ANALYSIS
\chapter{Experimental Results and Analysis}

\section{Results Summary}

The system was functionally evaluated using sample resume inputs and multiple
target role selections including Frontend Developer, Backend Developer, Full
Stack Developer, Data Scientist, DevOps Engineer, Android Developer, and custom
roles. The evaluation focused on three criteria: \textbf{Input Handling},
\textbf{Retrieval Relevance}, and \textbf{Output Completeness}.

\begin{table}[H]
\centering
\caption{Summary of functional evaluation results.}
\begin{tabular}{|p{4cm}|p{3cm}|p{6cm}|}
\hline
\textbf{Test Area} & \textbf{Status} & \textbf{Observation} \\
\hline
Plain text resume input & Successful & Resume text was accepted, validated, and analysed. \\
\hline
PDF resume upload & Successful & Text-based PDFs were parsed using \texttt{pypdf}. \\
\hline
Predefined role selection & Successful & Roadmap chunks were retrieved for selected roles. \\
\hline
Custom role input & Successful & Gemini inferred role-specific context with retrieved support chunks. \\
\hline
RAG retrieval & Successful & Relevant roadmap chunks were selected using embedding similarity. \\
\hline
JSON result rendering & Successful & Match score, gaps, strengths, and roadmap were rendered in the UI. \\
\hline
\end{tabular}
\end{table}

The system successfully generated structured gap analysis reports and learning
roadmaps for the tested role categories. The first RAG run creates the local
vector store, while later runs reuse the stored embeddings unless the embedding
model or roadmap count changes.

\section{Interpretation of Results}

\begin{itemize}
  \item \textbf{Gap Analysis Quality:} RAG-based context retrieval helps Gemini
        focus on role-relevant skill categories rather than relying only on a
        broad prompt.
  \item \textbf{Roadmap Usefulness:} The generated roadmaps provide phased skill
        development suggestions with resources, making the output actionable for
        candidates.
  \item \textbf{Response Time:} The first analysis can take longer because the
        roadmap vector store is created. Subsequent analyses are faster because
        stored vectors are reused.
  \item \textbf{Input Flexibility:} Supporting both pasted text and PDF upload
        improves usability for candidates with different resume formats.
\end{itemize}

\section{Sample Output: Gap Analysis Report}

Figure~5.1 shows a sample gap analysis report generated by the system for a
candidate applying for a Data Science or Machine Learning related role.

\begin{figure}[H]
  \centering
  % \includegraphics[width=0.88\textwidth]{sample_output.png}
  \fbox{\parbox{0.85\textwidth}{\centering \vspace{3cm}
  [Screenshot: Sample Gap Analysis Output - Missing Skills: MLOps,
  Kubernetes, Feature Stores | Roadmap: Week 1--2: MLOps fundamentals
  (Coursera), Week 3--4: Docker \& Kubernetes (KodeKloud), ...]
  \vspace{3cm}}}
  \caption{Sample gap analysis and roadmap output for a Data Science candidate.}
\end{figure}

\section{Comparison with Existing Approaches}

\begin{table}[H]
\centering
\caption{Comparison of the proposed system with existing approaches.}
\begin{tabular}{|p{3.5cm}|p{3cm}|p{3cm}|p{3cm}|}
\hline
\textbf{Approach} & \textbf{Method} & \textbf{Roadmap Support} & \textbf{Context Awareness} \\
\hline
Keyword Matching & Exact word overlap & No & Low \\
\hline
Resume Parser APIs & NLP + rules & No & Medium \\
\hline
Pure LLM & Prompt only & Yes & Medium \\
\hline
\textbf{Proposed System} & \textbf{RAG + Gemini + Flask} & \textbf{Yes} & \textbf{High} \\
\hline
\end{tabular}
\end{table}

The proposed RAG-based approach is more context-aware than keyword matching
because it retrieves semantically relevant roadmap chunks before generating the
final analysis.

% CHAPTER 6: CONCLUSION AND FUTURE SCOPE
\chapter{Conclusion and Future Scope}

\section{Conclusion}

This project has successfully demonstrated the development of an intelligent
Resume Gap Analysis system using Retrieval-Augmented Generation (RAG), Gemini,
LangChain, and Flask. The system addresses a real-world problem in career
development by comparing a candidate's resume with a target job role and
generating both a detailed gap analysis and a structured learning roadmap.

\begin{itemize}
  \item The RAG pipeline retrieves relevant roadmap skill chunks from a local
        SQLite vector store before generating the final analysis.
  \item Gemini embeddings enable semantic matching between the resume-role query
        and stored roadmap chunks.
  \item Gemini chat generation produces a structured report containing match
        score, summary, strengths, matched skills, missing skills, and roadmap.
  \item The Flask-based web interface provides an accessible, user-friendly
        platform for candidates and career guidance teams.
  \item The project keeps the API key on the backend and uses plain HTML and CSS
        for a lightweight frontend.
\end{itemize}

The project demonstrates that combining modern NLP techniques with a lightweight
web framework can produce practically useful AI applications for education,
placement preparation, and career guidance.

\section{Future Scope}

The following enhancements are planned for future iterations of the project:

\begin{enumerate}
  \item \textbf{OCR for Scanned Resumes:} Add OCR support for image-only PDF
        resumes using tools such as Tesseract.
  \item \textbf{Expanded Knowledge Base:} Add company-specific role requirements,
        course catalogs, certification details, and live learning resources.
  \item \textbf{Dedicated Vector Database:} Replace the local SQLite vector store
        with ChromaDB, FAISS, Pinecone, or Weaviate for larger-scale deployment.
  \item \textbf{Multi-Resume Batch Processing:} Enable HR teams to upload and
        analyse multiple resumes for a single target role.
  \item \textbf{Dynamic Roadmap Tracking:} Integrate user accounts where
        candidates can mark milestones as complete and receive updated roadmaps.
  \item \textbf{LinkedIn and GitHub Integration:} Supplement resumes with live
        profile and project evidence.
  \item \textbf{Multilingual Support:} Extend the system to support resumes in
        multiple languages using multilingual embedding models.
  \item \textbf{Streaming Responses:} Implement Server-Sent Events (SSE) to show
        LLM output progressively.
\end{enumerate}

% REFERENCES
\chapter*{References}
\addcontentsline{toc}{chapter}{References}

\begin{enumerate}

\item Lewis, P., Perez, E., Piktus, A., et al.\ (2020). Retrieval-Augmented
Generation for Knowledge-Intensive NLP Tasks. \textit{Advances in Neural
Information Processing Systems (NeurIPS)}, 33, 9459--9474.\\
\textgreater\space DOI: 10.48550/arXiv.2005.11401\\
\textit{(Foundational paper on RAG; used to justify the retrieval-augmented
approach for context-grounded generation.)}

\item Google AI for Developers.\ (2025). Gemini API Embeddings Documentation.\\
\textgreater\space \url{https://ai.google.dev/gemini-api/docs/embeddings}\\
\textit{(Reference for Gemini embedding models used to create semantic vectors
for roadmap retrieval.)}

\item Google AI for Developers.\ (2025). Gemini API Documentation.\\
\textgreater\space \url{https://ai.google.dev/gemini-api/docs}\\
\textit{(Reference for Gemini model usage in content generation.)}

\item Chase, H.\ (2022). LangChain: Building Applications with LLMs through
Composability. \textit{GitHub Repository.}\\
\textgreater\space \url{https://github.com/langchain-ai/langchain}\\
\textit{(LangChain framework used for prompt, model, and parser orchestration.)}

\item Grinberg, M.\ (2018). \textit{Flask Web Development: Developing Web
Applications with Python} (2nd ed.). O'Reilly Media.\\
\textit{(Reference for Flask application design and best practices.)}

\item Python Software Foundation.\ (n.d.). SQLite3 Python Documentation.\\
\textgreater\space \url{https://docs.python.org/3/library/sqlite3.html}\\
\textit{(Reference for the local SQLite vector store used in the project.)}

\item pypdf Contributors.\ (n.d.). pypdf Documentation.\\
\textgreater\space \url{https://pypdf.readthedocs.io/}\\
\textit{(Reference for PDF text extraction from uploaded resumes.)}

\end{enumerate}

\end{document}
